	\subsection*{Définition de l'indexation documentaire}
	
	\begin{itemize}
		\item \textbf{Indexation documentaire (Sumominen, Husson, Reignier)}
		
	définition de l'indexation : 
	"Le fait de dresser un répertoire ou une liste, généralement alphabétique, des sujets traités, des noms cités dans un ouvrage, suivis des références aux pages, aux paragraphes" (Rolland-Coul. 1969).
		\footcite{zotero-item-572}
		
	autre définition : 
	 Selon l'AFNOR, l'indexation est la ”représentation par les éléments d’un langage documentaire ou naturel, des
	notions résultant de l’analyse du contenu d’un document en vue d’en faciliter la recherche”
	\footcite{husson2025}
	
	«Les bibliothèques gèrent un grand nombre de métadonnées, pour différents types de documents. Le plus souvent, ces documents sont indexés à l'aide de mots sujets sélectionnés au sein d'un vocabulaire contrôlé, afin d'être recherchés plus tard. L'indexation manuelle de documents est un travail intellectuel très consommateur de temps.»
	\footcite[p.~21]{suominen2019a}
	
	définition de l'indexation :
	"L'indexation est cette technique consistant à caractériser le contenu d'un document et l'information qu'il détient de manière à le retrouver quand on effectue des recherches sur l'un des sujets dont il traite. La difficulté est donc de savoir caractériser et représenter l'information documentaire pour qu'il soit aisé de la mettre en rapport avec des sujets d'investigation. Mise en rapport d'une requête et d'un contenu représenté et synthétisé, l'indexation permet de s'orienter dans la masse des documents et d'organiser ses connaissances. L'indexation appartient donc à ce qu'on appelle depuis quelques années les techniques intellectuelles Indexation et archivage de contenus multimédias"
	\footcite{zotero-item-582}
	
	contrairement aux documents textuels qui comportent souvent des éléments d'indexation dans leur forme (couverture de livre, ) les archives audiovisuelles n'en contiennent peu ou pas, ce qui les rend opaques et peu récéptives aux pratiques d'indexation. Il s'agit alors de mettre en place une « l'indexation par le contenu ». \footcite{zotero-item-582}
	
	
	\item \textbf{Indexation automatique}
	
	«Les systèmes d'indexation automatique suivent généralement un processus dans lequel les documents textuels sont, dans un premier temps, "préparés" [...] les documents sont convertis dans une représentation vectorielle de leur fréquence d'apparition, ce qui, à l'aide d'un algorithme, permet d'établir des comparaisons entre les mots-sujets utilisés.»
	\footcite[p.~21]{suominen2019a}
	
	«Les algorithmes d'indexation automatique se répartissent selon deux types d'approche : lexicale et associative.»
	\footcite[p.~21]{suominen2019a}
	
	Pour les archives audiovisuelles, des processus techniques s'ajoutent : extraction d'entités depuis l'images, extraction d'entités depuis le son. La forme du médium cinématographique crée une opacité et une complexité pour l'instauration de pipeline algorithmiques.
	
	\item \textbf{Indexation iconographique}
	
	«L'indexation iconographique consiste à décrire le contenu des images à l'aide de descripteurs afin de faciliter l'accès aux œuvres, sans cependant dispenser sa consultation.»
	\footcite[p.~3]{husson2025}
	
	
	\item \textbf{Archives audiovisuelles (Bassaïsteguy (INA), Rattinger, Le Ber, Attard)}
	
	\item \textbf{IA et vision par ordinateur (Arnab (ViViT), Gong (AST), Trojahn, Yuan)}
	
	\item \textbf{Pipeline}
	
	Mise en place d'une pipeline IA : 
	Ce sont plusieurs outils combinés qui créent l'indexation : analyse audio avec l'ASR et la segmentation, analyse de l'image avec la computer vision et shot detection, le NLP avec le NER et term extraction. Cette pipeline permet d'extraire des métadonnées en lien avec des vocabulaires contrôlés et ainsi d'indexer les images. 
	«Each tool specializes in a single task, and is brought together with other tools in pipelines. Technologies can be layered upon one another in order to more meaningfully annotate content.»
	\footcite{2021}
	
	Implémentation d'une combinaison d'outils différents pour l'indexation : 
	"La combinaison de différents outils et leur intégration dans des environnementslogiciels existants est un sujet important, pour lequel il existe désormais un large éventail de solutions, souvent open source"
	\footcite{zotero-item-584}
	
	\item Outils
	
	Trint : automatic speech recognition 
	FasterR CNN et Farec CNN : face detection models 
	YOLO : outil opensource de détection d'objets et d'identification
	\item Cas d'indexation

	The vanderbilt university television news archive : ils utilisent Trint pour transcrire les discours et engagent des étudiants pour corriger les transcriptions sur l'interface de l'outil Trint : 
	
	"students also note the start and stop time of segments, and write relevant on-screen information such as names into the transcript. Each 30-minute program takes one hour of student worker time to correct. When the transcript is complete, it is exported to Amazon’s Named Entity Recognition service on the cloud, where information such as person, location, or topic is pulled out. That information is then brought back into the workflow with scripts written in Python that turn the information into individual titles for each news segment. Additional information such as reporter or location are added to database fields, and duplicate records are kept in PBCore for future use when more of the information from the Named Entity Recognition system might be stored in the database."
	\footcite{2021b}
	
	\end{itemize}

\subsection{limites}
«How to give context to the machines? To give common sense? It's a problem of artificial intelligence from the beginning.»
\footcite{2021}
	


